\documentclass[11pt]{article}

\usepackage[T1]{fontenc}
\usepackage[margin=1.1in]{geometry}
\usepackage{amsmath,amssymb,amsthm}
\usepackage{enumitem}
\usepackage{booktabs}
\usepackage[colorlinks=true,linkcolor=blue,citecolor=blue,urlcolor=blue]{hyperref}

\newtheorem{theorem}{Theorem}[section]
\newtheorem{proposition}[theorem]{Proposition}
\newtheorem{lemma}[theorem]{Lemma}
\newtheorem{corollary}[theorem]{Corollary}
\newtheorem{definition}[theorem]{Definition}
\newtheorem{conjecture}[theorem]{Conjecture}
\newtheorem{question}[theorem]{Question}
\theoremstyle{remark}
\newtheorem{remark}[theorem]{Remark}
\newtheorem{example}[theorem]{Example}

\newcommand{\R}{\mathbb{R}}
\newcommand{\X}{X}
\newcommand{\Xhat}{\widehat{\X}}
\newcommand{\Mx}{M_{\X}}
\newcommand{\Rx}{\mathcal{R}_{\X}}
\newcommand{\Hb}[1]{\mathcal{H}_{#1}}
\newcommand{\Sb}[1]{S_{#1}}
\newcommand{\dist}{\operatorname{dist}}
\newcommand{\conv}{\operatorname{conv}}
\newcommand{\cconv}[1]{\overline{\conv}\,#1}
\newcommand{\Cut}{\operatorname{Cut}}
\newcommand{\B}{B}
\newcommand{\bdry}{\partial_{\infty}}
\newcommand{\sca}[2]{\left\langle #1,\, #2\right\rangle}

\title{Split horoballs and the necessity direction of medial-axis
reconstruction on Hadamard manifolds\thanks{Companion note to
\cite{R1}, which generalised the sufficiency half of Theorem~6 of
\cite{Bia} to Riemannian manifolds, and to our formalization notes
\cite{Notes} on the Euclidean case. Drafted with the assistance of
Claude (Anthropic); all arguments were re-derived and checked by hand.
Remaining errors are ours.}}
\author{}
\date{June 2026}

\begin{document}
\maketitle

\begin{abstract}
Conjecture~5.1 of \cite{R1} proposed that on a Hadamard manifold every
reconstructible point of $\X^{c}$ lies in the horoball hull $\Xhat$ or
in a \emph{defective} maximal free horoball. The hull-based notion of
defectiveness is inadequate: the horoball hull can collapse to $\X$
itself (already for a two-point set in $\mathbb{H}^{2}$), destroying
all defectiveness while reconstruction survives. We replace
defectiveness by tangency data on the horosphere --- ultimately the
\emph{split} condition of \S\ref{sec:sufficiency} --- and prove the
structural half of the corresponding necessity statement. The key
negative result (Proposition~\ref{prop:rigidity}) is a rigidity
statement: inflating a medial ball along the geodesic through one of
its feet is cut \emph{instantly} by any other foot, so escaping
inflations exist only from uniquely-footed centres and produce
horoballs with a single attained contact --- bitangency can only be
created by equidistant motion of centres. The positive results: a
drift theorem (Theorem~\ref{thm:branch-two}) --- medial balls of
unbounded radius containing $p$ with uniform margin converge, via the
horofunction boundary, to a maximal free horoball containing $p$ at
definite depth, whose attained contacts are the subsequential limits
of feet --- and an unconditional structure trichotomy
(Theorem~\ref{thm:structure}): \emph{every} point of $\X^{c}$ lies in
$\Xhat$, or satisfies the drift hypothesis, or lies at definite depth
in single-contact escape horoballs; there is no residual case. In
$\R^{n}$, reconstructible points outside the hull always drift
(Proposition~\ref{prop:euclidean}, by lifting the defect witness of
\cite{Notes}). The arguments use no entry points into $\Xhat$ and are
immune to the tangential-entry gap of \cite[\S5]{R1}. We close with the taxonomy separating bitangency from
defectiveness, an example (a horocyclic arc in $\mathbb{H}^{2}$)
showing that bitangency alone does \emph{not} imply reconstruction,
the corrected notion of a \emph{split} horoball (separated contact
clusters), a complete proof that split horoballs are reconstructed
whenever the splitting scale exceeds half the asymptotic separation of
the contact rays --- unconditionally under strictly negative
curvature ---, a deepest-direction dichotomy reducing the remaining
necessity gap to an asymptotic-tangency regime along the deepest
axis, its resolution on $\mathbb{H}^{2}$ for non-degenerate drift, an
asymptotic-strand example showing that defectiveness survives exactly
where splitness fails. Finally (\S\ref{sec:dim3}) we show that the
split-or-defective conjecture is \emph{false} for $n\ge3$: in
$\mathbb{H}^{3}$ the unit circle on a horosphere has collapsing hull
and only connected (hence non-split, non-defective) contact sets, yet
$\Rx=\{z>1\}\cup\mathcal{H}_{0}\neq\varnothing$. Splitness is a
two-dimensional artefact: the operative condition is failure of the
contact set to be convex inside the intrinsically flat horosphere
$\cong\mathbb{E}^{n-1}$, which on a line ($n=2$) means a gap (split)
but for $n\ge3$ is implied already by a connected set such as a
circle. We propose the corrected notion --- horospherical defect --- and
prove its sufficiency direction in $\mathbb{H}^{n}$ via a duality with
Motzkin's theorem (a closed Euclidean set is convex iff Chebyshev):
every horospherically defective horoball with attained, bulk-isolated
contacts satisfies $\Hb{}\subseteq\Rx$, and when $\X$ lies on the
horosphere this is sharp ($\Rx=\varnothing$ for convex $\X$). In
particular the circle counterexample is fully resolved ---
$\Rx=\{z>1\}\cup\mathcal{H}_{0}$ is exactly the union of its
horospherically defective horoballs --- so the example that refuted
the old conjecture confirms its repair. For necessity we prove the non-merging estimate: the two feet of a
reconstructing drift, once they ascend to the horosphere
(Lemma~\ref{lem:ascent}), either collapse onto a single contact or
stay separated by exactly
$\operatorname{arccosh}(1+\tfrac12\lvert\kappa-\kappa'\rvert^{2})$,
with no intermediate rate; in the non-collapse case their distinct
limits are equidistant-nearest contacts, so by Motzkin $K$ is
non-convex and the horoball is horospherically defective
(Theorem~\ref{thm:nonmerge}). This proves necessity in $\mathbb{H}^{n}$
for every bounded, non-collapsing drift, reducing the whole conjecture
to a single residual asymptotic/local-hull question. A status table
closes the note.
\end{abstract}

\section{Setting and definitions}\label{sec:setting}

Throughout, $(M^{n},g)$ is a \emph{Hadamard manifold}: complete,
simply connected, of nonpositive sectional curvature. We write
$d(y,z)$ for the Riemannian distance. Every geodesic is globally
minimizing, extends to a complete geodesic line, and is determined by
a point and a velocity. $\X\subsetneq M$ is a fixed nonempty closed
set, $d(y):=\dist(y,\X)$, and
\[
m(y):=\{x\in\X \mid d(y,x)=d(y)\}
\]
is the nonempty compact set of \emph{feet} of $y$ (Hopf--Rinow). All
balls $\B(y,r)$ are open. A ball (or any set) is \emph{free} if it is
disjoint from $\X$; thus $\B(y,r)$ is free iff $r\le d(y)$.

\begin{definition}[Medial axis, reconstructible set]\label{def:medial}
The \emph{medial axis} is
$\Mx:=\{a\in\X^{c} : \#m(a)\ge 2\}$; on a Hadamard manifold this
coincides with the set of points admitting at least two distinct
minimizing geodesics to $\X$ (no branching, cf.\
\cite[Rem.~2.2]{R1}). The \emph{reconstructible set} is
\[
\Rx:=\bigcup_{a\in\Mx}\B\bigl(a,d(a)\bigr).
\]
For $p\in\Rx$ and $a\in\Mx$ with $p\in\B(a,d(a))$ we call
$\mu(a,p):=d(a)-d(p,a)>0$ the \emph{margin} of $p$ in the medial ball
$\B(a,d(a))$.
\end{definition}

\begin{definition}[Busemann function, maximal free horoball]
\label{def:busemann}
For a unit-speed geodesic ray $\gamma:[0,\infty)\to M$ the
\emph{Busemann function} is
$b_{\gamma}(y):=\lim_{t\to\infty}\bigl(d(y,\gamma(t))-t\bigr)$; the
limit exists since the expression is non-increasing in $t$ and bounded
below by $-d(y,\gamma(0))$. It is $1$-Lipschitz, convex, and
$b_{\gamma}(\gamma(s))=-s$. Set $c_{\gamma}:=\inf_{\X}b_{\gamma}$.
When $c_{\gamma}$ is finite, the \emph{maximal free horoball} and its
\emph{horosphere} are
\[
\Hb{\gamma}:=\{b_{\gamma}<c_{\gamma}\},\qquad
\Sb{\gamma}:=\{b_{\gamma}=c_{\gamma}\}.
\]
$\Hb{\gamma}$ is open, convex, free, and maximal among free
sub-level sets of $b_{\gamma}$. The \emph{depth} of a point
$p\in\Hb{\gamma}$ is $c_{\gamma}-b_{\gamma}(p)>0$. The
\emph{horoball hull} is
$\Xhat:=M\setminus\bigcup\{\text{free open horoballs}\}
=\{y : b_{\gamma}(y)\ge c_{\gamma}\ \text{for every ray}\
\gamma\}$, a closed set containing $\X$.
\end{definition}

\begin{definition}[Bitangent; defective]\label{def:bitangent}
A maximal free horoball $\Hb{\gamma}$ is
\begin{enumerate}[label=(\alph*)]
\item \emph{bitangent} if $\Sb{\gamma}\cap\X$ contains at least two
points (each contact point is the foot of a distinct minimizing
geodesic from any centre $\gamma(t)$ deep enough, so bitangency
witnesses genuine multiplicity on Hadamard manifolds);
\item \emph{defective} (the notion of \cite[Def.~2.5]{R1}) if
$(\Sb{\gamma}\cap\Xhat)\setminus\X\neq\emptyset$.
\end{enumerate}
\end{definition}

\begin{remark}[Why defectiveness must be replaced]
\label{rem:collapse}
For $\X=\{x_{1},x_{2}\}\subset\mathbb{H}^{2}$ one has $\Xhat=\X$: in
the upper half-plane model with $x_{1,2}=(\pm s,1)$, every point
$z=(u,v)\notin\X$ with $v\le 1$ lies in the free horoball tangent to
the ideal boundary at $(u,0)$ of Euclidean radius
$\rho=\tfrac12\bigl(\min_{\pm}(u\mp s)^{2}+1\bigr)$, and every point
with $v>1$ lies in the free horoball $\{y>1\}$. (This corrects
Remark~2.6 of \cite{R1}: the horoball hull does not even contain the
geodesic segment $[x_{1},x_{2}]$, whose interior dips into
$\{y>1\}$.) Consequently no maximal free horoball is defective,
although $\Rx=\{y>1\}\cup\B_{0}\neq\emptyset$, where $\B_{0}$ is the
horoball tangent at $0$ through both points: Conjecture~5.1 of
\cite{R1} fails as stated. Both horoballs constituting $\Rx$ are
bitangent, which motivates Definition~\ref{def:bitangent}(a) and the
revision below. We do not reproduce the full computation here, the
original conjecture having been a working formulation.
\end{remark}

\begin{conjecture}[Revised necessity-and-sufficiency]
\label{conj:revised}
Let $M$ be Hadamard and $\X$ closed. Then
\[
\Rx \;=\; (\Xhat\setminus\X)\;\cup\;
\bigcup\Bigl\{\Hb{\gamma} \;:\; \Hb{\gamma}\ \text{maximal free,
split (Def.~\ref{def:split}) or defective}\Bigr\}.
\]
\end{conjecture}

The inclusion $\supseteq$ holds for $\Xhat\setminus\X$ by
\cite[Thm.~3.1]{R1} (completeness only) and for defective horoballs by
\cite[Thm.~4.1]{R1}; for split horoballs it is proved under a
separation-scale condition in Theorem~\ref{thm:split} below.
Bitangency alone is \emph{not} sufficient (Example~\ref{ex:arc}),
which is what forces the split refinement. The present note proves the structural half of
$\subseteq$ (Proposition~\ref{prop:rigidity},
Theorems~\ref{thm:branch-two} and \ref{thm:structure}) and the split
half of $\supseteq$ (Theorem~\ref{thm:split}). One external
input is used, as in \cite{R1}:

\begin{lemma}[Cut-locus density; Wolter, Albano \cite{Wol,Alb}]
\label{lem:cut}
Let $M$ be complete and $\X$ closed. The cut locus $\Cut(\X)$ --- the
set of endpoints of maximal minimizing segments of normal geodesics
from $\X$ --- satisfies $\Cut(\X)\subseteq\overline{\Mx}$.
\end{lemma}

\section{The inflation dichotomy, with margins}\label{sec:inflation}

We first record the inflation construction of \cite[Thm.~3.1]{R1} in
the quantitative form we need. Fix $a\in\X^{c}$, a foot $q\in m(a)$,
and let $\gamma:[0,\infty)\to M$ be the unit-speed ray with
$\gamma(0)=q$ and $\gamma(\rho)=a$, where $\rho:=d(a)$ (the geodesic
from $q$ through $a$, extended; on a Hadamard manifold the extension
is a globally minimizing ray). Define the set of \emph{free times}
$T:=\{t>0 : \B(\gamma(t),t)\cap\X=\emptyset\}$ and the
\emph{escape time}
\[
r(a,q):=\sup T\;\in\;[\rho,\infty].
\]

\begin{lemma}[Monotone free inflation]\label{lem:monotone}
$T$ is an interval containing $(0,\rho]$, and for $t'<t$ with
$t\in T$ one has $\B(\gamma(t'),t')\subseteq\B(\gamma(t),t)$.
Moreover for every $t\in T$:
\begin{enumerate}[label=(\roman*)]
\item $d(\gamma(t))=t$, with $q\in m(\gamma(t))$;
\item if $p\in\B(a,\rho)$ with margin $\mu=\rho-d(p,a)$, then
$d(p,\gamma(t))\le t-\mu$ for all $t\ge\rho$, $t \in T$;
\item every $q'\in m(a)$ satisfies $d(q',\gamma(t))=t$; in
particular all feet of $a$ lie on $\partial\B(\gamma(t),t)$.
\end{enumerate}
\end{lemma}

\begin{proof}
Since $d(\gamma(t'),\gamma(t))=t-t'$ (geodesics minimize), for
$y\in\B(\gamma(t'),t')$ we get
$d(y,\gamma(t))\le d(y,\gamma(t'))+(t-t')<t$, proving the nesting and
that $T$ is an interval; $\rho\in T$ because $\B(a,\rho)$ is free.
(i): freeness gives $d(\gamma(t))\ge t$, while
$d(\gamma(t),q)=t$ gives $\le$ and exhibits $q$ as a foot. (ii):
$d(p,\gamma(t))\le d(p,a)+d(a,\gamma(t))=(\rho-\mu)+(t-\rho)$. (iii):
$d(q',\gamma(t))\le d(q',a)+(t-\rho)=\rho+(t-\rho)=t$ by the triangle
inequality, and $\ge t$ by freeness, since $q'\in\X$.
\end{proof}

Part (iii) looks like an engine for producing bitangent limit
horoballs; in fact it is self-limiting. Boundary-riding is an
\emph{equality} in the triangle inequality, and equalities in uniquely
geodesic spaces are rigid: they force collinearity, which destroys
freeness immediately. Section~\ref{sec:branch-one} turns this into a
rigidity statement that any necessity mechanism must respect.

\begin{lemma}[Cut branch]\label{lem:cutbranch}
Suppose $r:=r(a,q)<\infty$. Then $c:=\gamma(r)$ satisfies $d(c)=r$,
the ball $\B(c,r)$ is free and contains $\B(a,\rho)$, and
$c\in\Cut(\X)\subseteq\overline{\Mx}$. Consequently, if
$p\in\B(a,\rho)$ has margin $\mu$, then for every
$\varepsilon\in(0,\mu/2)$ there exists $a'\in\Mx$ with
\[
d(a')\ge r-\varepsilon,\qquad
p\in\B\bigl(a',d(a')\bigr),\qquad
\mu(a',p)\ge\mu-2\varepsilon .
\]
\end{lemma}

\begin{proof}
$d(\gamma(t))=t$ on $T$ and continuity of $d$ give $d(c)=r$. If
$y\in\B(c,r)$, pick $\eta>0$ with $d(y,c)<r-\eta$; for $t<r$ with
$r-t<\eta/2$,
$d(y,\gamma(t))\le d(y,c)+(r-t)<r-\eta+\eta/2<t$, so
$y\in\B(\gamma(t),t)$, which is free; hence $\B(c,r)$ is free, and it
contains each $\B(\gamma(t),t)\supseteq\B(a,\rho)$ by
Lemma~\ref{lem:monotone}. The segment $\gamma|_{[0,t]}$ minimizes
$d(\cdot,\X)$ for every $t\le r$ by Lemma~\ref{lem:monotone}(i), while
for $t>r$ freeness fails, so $d(\gamma(t))<t$ and the segment no
longer minimizes: $c$ is the cut point of $\X$ along $\gamma$, and
Lemma~\ref{lem:cut} applies. Finally take $a'\in\Mx$ with
$d(a',c)<\varepsilon$: then $d(a')\ge d(c)-\varepsilon=r-\varepsilon$
and
$d(p,a')\le d(p,c)+\varepsilon\le(r-\mu)+\varepsilon$ using
Lemma~\ref{lem:monotone}(ii) at $t\to r$; hence
$\mu(a',p)=d(a')-d(p,a')\ge\mu-2\varepsilon>0$.
\end{proof}

\section{Rigidity: no escape from a medial centre}
\label{sec:branch-one}

\begin{proposition}[Instant cutting]\label{prop:rigidity}
Let $M$ be Hadamard and $a\in\Mx$. Then $r(a,q)=d(a)$ for every foot
$q\in m(a)$: the inflation of a medial ball along any of its
foot-geodesics is cut instantly.
\end{proposition}

\begin{proof}
Write $\rho:=d(a)$, pick a second foot $q'\neq q$, and suppose
$t\in T$ with $t>\rho$. By Lemma~\ref{lem:monotone}(iii),
$d(q',\gamma(t))=t=d(q',a)+d(a,\gamma(t))$: equality in the triangle
inequality, so the concatenation of the geodesic segments $[q',a]$
and $[a,\gamma(t)]$ has length $d(q',\gamma(t))$ and is therefore the
(unique) geodesic from $q'$ to $\gamma(t)$. In particular it passes
through $a$ with velocity $\dot\gamma(\rho)$, so its initial segment
$[q',a]$, reversed, is part of $\gamma$:
$q'=\gamma(\rho-d(q',a))=\gamma(0)=q$, contradicting $q'\neq q$.
Hence $T\subseteq(0,\rho]$ and $r(a,q)=\rho$.
\end{proof}

\begin{remark}\label{rem:vacuity}
Proposition~\ref{prop:rigidity} redirects the search for a necessity
mechanism. Escaping inflations cannot start at medial centres, so a
single inflation can never tag two feet on a limit horosphere;
bitangency must instead be produced by moving centres \emph{along
equidistant loci} of two contact clusters --- which is exactly what
the pincer in the proof of Theorem~\ref{thm:split} does, and what the
medial families of the worked examples (e.g.\ the bisector family of
a two-point set in $\mathbb{H}^{2}$) actually look like. An earlier
draft of this note stated a bitangency theorem for escaping
inflations from medial centres; the rigidity above shows its
hypothesis is unsatisfiable.
\end{remark}

What survives of the escape mechanism is its single-contact version,
which by Proposition~\ref{prop:rigidity} can only fire at
uniquely-footed centres; it is used in the proof of the structure
theorem below.

\begin{theorem}[Single-foot escape]\label{thm:branch-one}
Let $M$ be Hadamard, $\X$ closed, $a\in\X^{c}$, $q\in m(a)$, and let
$p\in\B(a,d(a))$ have margin $\mu>0$. If $r(a,q)=\infty$ --- which
forces $m(a)=\{q\}$ --- then with $\gamma$ the ray of
\S\ref{sec:inflation}:
\begin{enumerate}[label=(\alph*)]
\item $c_{\gamma}=0$, attained at $q$: $\Hb{\gamma}=\{b_{\gamma}<0\}$
is a maximal free horoball with attained contact
$q\in \Sb{\gamma}\cap\X$;
\item $b_{\gamma}(p)\le-\mu$, so $p\in\Hb{\gamma}$ at depth at least
$\mu$;
\item $\Hb{\gamma}=\bigcup_{t>0}\B(\gamma(t),t)$, an increasing union
of free balls each containing $\B(a,d(a))$ for $t\ge d(a)$.
\end{enumerate}
\end{theorem}

\begin{proof}
All balls $\B(\gamma(t),t)$, $t>0$, are free. For $x\in\X$ this
means $d(x,\gamma(t))\ge t$, hence $b_{\gamma}(x)\ge0$ and
$c_{\gamma}\ge0$. By Lemma~\ref{lem:monotone}(i),
$d(q,\gamma(t))=t$ for all $t$, so $b_{\gamma}(q)=0$; with $q\in\X$
this forces $c_{\gamma}=0$, attained at $q$, proving (a). By
Lemma~\ref{lem:monotone}(ii), $d(p,\gamma(t))\le t-\mu$ for all
$t\ge d(a)$, whence $b_{\gamma}(p)\le-\mu<0=c_{\gamma}$, proving (b).
For (c), the inclusion $\subseteq$ of the union follows from
$b_{\gamma}(y)\le d(y,\gamma(t))-t<0$ for $y\in\B(\gamma(t),t)$;
conversely if $b_{\gamma}(y)<0$ then $d(y,\gamma(t))<t$ for all large
$t$. The containment statement is Lemma~\ref{lem:monotone}.
\end{proof}

\section{Branch II: unbounded medial radii}\label{sec:branch-two}

By Proposition~\ref{prop:rigidity}, escaping inflations cannot start
from medial centres; the mechanism by which the medial axis actually
produces horoballs is \emph{drift}: medial balls containing $p$ with
radii tending to infinity and uniform margin (realized, for instance,
by the bisector family of a two-point set in $\mathbb{H}^{2}$). For
this branch we use the horofunction description of the boundary at
infinity. For $\delta>0$ set
\[
T_{\delta}(p):=\sup\bigl\{\,d(a) : a\in\Mx,\ \mu(a,p)\ge\delta\,\bigr\}
\in[0,\infty],\qquad \sup\emptyset:=0 .
\]

\begin{lemma}[Horofunction convergence; {\cite[II.8.13]{BH}}]
\label{lem:horofn}
Let $M$ be Hadamard, $o\in M$, and let $(a_{k})\subset M$ with
$d(o,a_{k})\to\infty$ and $a_{k}\to\xi\in\bdry M$ in the cone
topology. Then
$h_{k}(\cdot):=d(\cdot,a_{k})-d(o,a_{k})\longrightarrow
b_{\xi}(\cdot)$ uniformly on compact sets, where $b_{\xi}$ is the
Busemann function of the ray from $o$ to $\xi$, normalized by
$b_{\xi}(o)=0$. Moreover $\overline{M}=M\cup\bdry M$ is compact, so
every unbounded sequence subconverges to some $\xi\in\bdry M$.
\end{lemma}

\begin{theorem}[Limit horoball from unbounded medial radii]
\label{thm:branch-two}
Let $M$ be Hadamard, $\X$ closed, $p\in\X^{c}$, and suppose
$T_{\delta}(p)=\infty$ for some $\delta>0$. Then there exist
$\xi\in\bdry M$, $s_{\infty}\in[\delta,\,d(p)]$, medial centres
$a_{k}\in\Mx$ with $r_{k}:=d(a_{k})\to\infty$ and
$\mu(a_{k},p)\ge\delta$, such that, with $b_{\xi}$ normalized at
$o=p$:
\begin{enumerate}[label=(\alph*)]
\item $\X\subseteq\{b_{\xi}\ge s_{\infty}\}$ and
$b_{\xi}(p)=0\le s_{\infty}-\delta$; hence $c_{\xi}\ge
s_{\infty}>0=b_{\xi}(p)$ and $p$ lies in the maximal free horoball
$\Hb{\xi}=\{b_{\xi}<c_{\xi}\}$ at depth at least $\delta$;
\item the free balls exhaust the sub-horoball: every compact
$K\subseteq\{b_{\xi}<s_{\infty}\}$ satisfies
$K\subseteq\B(a_{k},r_{k})$ for all large $k$;
\item \emph{(contact capture)} if feet $q_{k}\in m(a_{k})$ have a
bounded subsequence, then along it $q_{k}\to q_{*}\in\X$ with
$b_{\xi}(q_{*})=s_{\infty}$; in that case $c_{\xi}=s_{\infty}$, the
infimum is attained at $q_{*}$, and $q_{*}\in \Sb{\xi}\cap\X$.
If two sequences of feet $q_{k}\neq q_{k}'$ of $a_{k}$ subconverge to
\emph{distinct} limits, $\Hb{\xi}$ is bitangent.
\end{enumerate}
\end{theorem}

\begin{proof}
Choose $a_{k}\in\Mx$ with $\mu(a_{k},p)\ge\delta$ and
$r_{k}=d(a_{k})\to\infty$. Set $t_{k}:=d(p,a_{k})\le r_{k}-\delta$
and $s_{k}:=r_{k}-t_{k}\ge\delta$. For any $x_{0}\in\X$, freeness of
$\B(a_{k},r_{k})$ gives $d(x_{0},a_{k})\ge r_{k}$, so
$d(x_{0},p)\ge d(x_{0},a_{k})-t_{k}\ge s_{k}$; hence
$s_{k}\le d(p)$, and after passing to a subsequence
$s_{k}\to s_{\infty}\in[\delta,d(p)]$. Also
$t_{k}\ge r_{k}-d(p)\to\infty$, so $(a_{k})$ is unbounded and, after a
further subsequence, $a_{k}\to\xi\in\bdry M$
(Lemma~\ref{lem:horofn}).

(a) For $x\in\X$: $h_{k}(x)=d(x,a_{k})-t_{k}\ge r_{k}-t_{k}=s_{k}$,
and letting $k\to\infty$, $b_{\xi}(x)\ge s_{\infty}$. Trivially
$b_{\xi}(p)=0$. Hence $c_{\xi}=\inf_{\X}b_{\xi}\ge
s_{\infty}\ge\delta$, while $c_{\xi}\le b_{\xi}(x_{0})<\infty$; the
maximal free horoball $\Hb{\xi}$ exists and contains $p$ at depth
$c_{\xi}-b_{\xi}(p)\ge\delta$.

(b) If $b_{\xi}(y)<s_{\infty}$ then, by locally uniform convergence
on a compact neighbourhood of $K$,
$h_{k}(y)<s_{k}$ for large $k$ uniformly on $K$, i.e.\
$d(y,a_{k})<t_{k}+s_{k}=r_{k}$.

(c) Along the bounded subsequence,
$h_{k}(q_{k})=d(q_{k},a_{k})-t_{k}=r_{k}-t_{k}=s_{k}$, and locally
uniform convergence plus $q_{k}\to q_{*}$ gives
$b_{\xi}(q_{*})=s_{\infty}$. Since $q_{*}\in\X$ (closedness) we get
$c_{\xi}\le s_{\infty}$, so $c_{\xi}=s_{\infty}$ is attained at
$q_{*}$. Two distinct such limits give two points of
$\Sb{\xi}\cap\X$.
\end{proof}

\begin{remark}\label{rem:escape-feet}
Without the boundedness hypothesis in (c) the feet may escape to
infinity and the contact be only asymptotic
($c_{\xi}=s_{\infty}$ not attained, or attained at fewer points than
the medial multiplicity suggests). This is a genuine phenomenon
already in $\R^{2}$: for
$\X=\operatorname{graph}(e^{x})\cup\{(5,0)\}$ the medial centres
$(\bar x_{T},-R_{T})$ with feet $(5,0)$ and
$(-T,e^{-T})$ have $R_{T}\to\infty$, one foot converging to the
attained contact $(5,0)$ and the other escaping along the asymptote;
the limit half-plane $\{y<0\}$ is maximal free with a single attained
contact. See \S\ref{sec:taxonomy}.
\end{remark}

\section{The structure theorem}\label{sec:structure}

\begin{theorem}[Structure trichotomy --- no residual case]
\label{thm:structure}
Let $M$ be Hadamard, $\X$ closed, and $p\in\X^{c}$ (reconstructible
or not). Then at least one of the following holds.
\begin{enumerate}[label=(\Roman*)]
\item \textbf{(Hull)} $p\in\Xhat\setminus\X$.
\item \textbf{(Drift)} $T_{\delta}(p)=\infty$ for some $\delta>0$;
then $p$ lies, at depth $\ge\delta$, in a maximal free horoball
exhausted (below level $s_{\infty}$) by medial balls, whose attained
contacts are the subsequential limits of their feet
(Theorem~\ref{thm:branch-two}).
\item \textbf{(Ray escape)} There is $\delta'>0$ such that for all
sufficiently large $t$, $p$ lies at depth $\ge\delta'$ in a maximal
free horoball with an attained contact, produced by an escaping
single-foot inflation (Theorem~\ref{thm:branch-one}).
\end{enumerate}
\end{theorem}

\begin{proof}
Assume (I) fails; since $p\notin\X$, this means $p\notin\Xhat$.
Choose a ray $\gamma$ with $b_{\gamma}(p)<c_{\gamma}=:c$ (finite by
this very inequality) and set $\delta^{*}:=c-b_{\gamma}(p)>0$. For
$x\in\X$ the function $s\mapsto d(x,\gamma(s))-s$ is non-increasing
with limit $b_{\gamma}(x)\ge c$, so $d(x,\gamma(t))\ge t+c$ for every
$t$: the balls $\B(\gamma(t),t+c)$ are free, and
$d(\gamma(t))\ge t+c$. Since $d(p,\gamma(t))-t\downarrow
b_{\gamma}(p)=c-\delta^{*}$, there is $t_{0}$ with
$d(p,\gamma(t))\le t+c-\delta^{*}/2$ for $t\ge t_{0}$: $p$ lies in
the maximal free ball $\B(\gamma(t),d(\gamma(t)))$ with margin
$\mu_{t}\ge\delta^{*}/2$. For $t\ge t_{0}$ pick a foot
$q_{t}\in m(\gamma(t))$. If $r(\gamma(t),q_{t})<\infty$ for an
unbounded set of $t$, Lemma~\ref{lem:cutbranch} with
$\varepsilon=\delta^{*}/8$ supplies $a'_{t}\in\Mx$ with
$d(a'_{t})\ge d(\gamma(t))-\delta^{*}/8\ge t+c-\delta^{*}/8$ and
$\mu(a'_{t},p)\ge\delta^{*}/4$, so $T_{\delta^{*}/4}(p)=\infty$:
case (II). Otherwise $r(\gamma(t),q_{t})=\infty$ for all large $t$,
and Theorem~\ref{thm:branch-one} (applied to $a=\gamma(t)$, margin
$\delta^{*}/2$) gives case (III) with $\delta'=\delta^{*}/2$.
\end{proof}

\begin{remark}
The trichotomy is unconditional: it removes the ``confinement'' case
conjectured in an earlier draft, and it does not require
$p\in\Rx$. Reconstructibility enters only through the open
\emph{split upgrade} (Question~\ref{q:merge}): for $p\in\Rx$, the
horoballs delivered by (II) and (III) must be improved to
split-or-defective ones. Thus Conjecture~\ref{conj:revised} reduces
to the split upgrade together with the sufficiency results of
\S\ref{sec:sufficiency}.
\end{remark}

In Euclidean space the drift branch always carries reconstruction
outside the hull:

\begin{proposition}[In $\R^{n}$, reconstruction outside the hull
forces drift]\label{prop:euclidean}
For $M=\R^{n}$ and $p\in\Rx\setminus\cconv\X$, case (II) of
Theorem~\ref{thm:structure} holds: $T_{\delta}(p)=\infty$ for some
$\delta>0$.
\end{proposition}

\begin{proof}
By the repaired necessity direction of Theorem~6 \cite[\S4.6]{Notes}
there are a unit vector $v$, $c=\sup_{\X}\sca{v}{\cdot}$, and a
defect witness
$w\in(\cconv\X\cap\{\sca{v}{\cdot}=c\})\setminus\X$ with
$h:=\sca{v}{p}-c>0$. Lift the \emph{witness}: $w_{t}:=w+tv$. Then
$d(w_{t})\ge\sca{v}{w_{t}}-c=t$ since
$\X\subseteq\{\sca{v}{\cdot}\le c\}$, while
$|p-w_{t}|^{2}=t^{2}-2th+|p-w|^{2}$ gives
$|p-w_{t}|\le t-h+O(1/t)$: for all large $t$, $p$ lies in
$\B(w_{t},d(w_{t}))$ with margin $\mu_{t}\ge h/2$.

No inflation from $w_{t}$ escapes. Indeed, suppose
$r(w_{t},q_{t})=\infty$ for a foot $q_{t}\in m(w_{t})$, and let
$u_{t}$ be the direction of the inflation ray, so
$q_{t}=w_{t}-\rho_{t}u_{t}$ with $\rho_{t}:=d(w_{t})$. Freeness of
all balls $\B(\gamma_{t}(s),s)$ gives, for every $x\in\X$,
$\sca{x-q_{t}}{u_{t}}\le0$; the same then holds for every point of
$\cconv\X$, in particular for $w$. Substituting
$w-q_{t}=\rho_{t}u_{t}-tv$:
\[
0\;\ge\;\sca{w-q_{t}}{u_{t}}\;=\;\rho_{t}-t\sca{v}{u_{t}}
\;\ge\;\rho_{t}-t\;\ge\;0 ,
\]
forcing equality throughout: $\sca{v}{u_{t}}=1$, hence $u_{t}=v$,
$\rho_{t}=t$, and $q_{t}=w_{t}-tv=w\notin\X$ --- contradicting
$q_{t}\in\X$.

Hence $r(w_{t},q_{t})<\infty$ for every $t$ large, and
Lemma~\ref{lem:cutbranch} with $\varepsilon=h/8$ supplies
$a'_{t}\in\Mx$ with $d(a'_{t})\ge d(w_{t})-h/8\ge t-h/8$ and
$\mu(a'_{t},p)\ge h/4$. Thus $T_{h/4}(p)=\infty$.
\end{proof}

\begin{corollary}[Strengthened Euclidean necessity]
\label{cor:euclidean}
For closed $\X\subseteq\R^{n}$, every $p\in\Rx\setminus\cconv\X$ lies
at positive depth in a maximal free half-space $L^{+}$ that is a limit
of medial balls in the sense of
Theorem~\ref{thm:branch-two}(b); the attained contacts of $\partial
L^{+}$ are the subsequential limits of their feet. In particular the
union in Theorem~6 of \cite{Bia} may be taken over maximal free
half-spaces arising as limits of medial balls.
\end{corollary}

\section{Bitangent versus defective}\label{sec:taxonomy}

The two notions are incomparable in general; they agree for compact
$\X$ in $\R^{n}$.

\begin{proposition}\label{prop:taxonomy}
\begin{enumerate}[label=(\alph*)]
\item \textbf{(Bitangent $\not\Rightarrow$ defective, Hadamard)} For
$\X$ a two-point set in $\mathbb{H}^{2}$
(Remark~\ref{rem:collapse}), the two bitangent maximal free horoballs
are not defective, since $\Xhat=\X$.
\item \textbf{(Defective $\not\Rightarrow$ bitangent, already in
$\R^{n}$)} For $\X=\operatorname{graph}(e^{x})\cup\{(5,0)\}
\subseteq\R^{2}$, the maximal free half-space $\{y<0\}$ has the single
attained contact $(5,0)$, yet it is defective: points $(3-\theta,0)$,
$\theta\ge0$ small, lie in $\cconv\X\setminus\X$ on the boundary line.
Mass escaping to infinity in the hull can even produce defective
half-spaces whose witness is unreachable by contact sequences: for
$\X=\{(5,0)\}\cup\{(-k^{2},-k):k\in\mathbb{N}\}$, the half-space
$\{y<-0\}$\,\footnote{Precisely: the maximal free half-space in
direction $-e_{2}$ through the contact $(5,0)$.} is defective with
witnesses $(5-c,0)$, $c\in(0,C]$, obtained from convex combinations
with weights $\sim1/k$ on $(-k^{2},-k)$, while every sequence in $\X$
with $y\to0$ is eventually constant equal to $(5,0)$.
\item \textbf{(Characterisation for compact $\X$ in $\R^{n}$)} If
$\X\subseteq\R^{n}$ is compact, a maximal free half-space
$L^{+}=\{\sca{v}{\cdot}>c\}$, $c=\max_{\X}\sca{v}{\cdot}$, is
defective if and only if the slice $\partial L^{+}\cap\X$ is
\emph{not convex}. In particular bitangency
($\#(\partial L^{+}\cap\X)\ge2$) is necessary but not sufficient for
defectiveness: a segment has a convex slice
(cf.\ Example~\ref{ex:arc}).
\end{enumerate}
\end{proposition}

\begin{proof}
(a) Remark~\ref{rem:collapse}. (b) The witness claims follow from
computing limits of convex combinations as indicated; closedness of
each $\X$ is immediate, and freeness and maximality of the half-spaces
follow since $\inf$ is $0$, attained at $(5,0)$ in both cases. (c)
($\Leftarrow$) If the slice is not convex, pick
$x_{1},x_{2}\in\partial L^{+}\cap\X$ and
$z\in[x_{1},x_{2}]\setminus\X$; then
$z\in\conv(\X\cap\partial L^{+})\setminus\X\subseteq
(\cconv\X\cap\partial L^{+})\setminus\X$ is a witness.
($\Rightarrow$) For compact
$\X$, $\conv\X$ is compact and
$\cconv\X\cap\partial L^{+}=\conv(\X\cap\partial L^{+})$: writing
$w=\sum\lambda_{i}x_{i}$ with $x_{i}\in\X$, $\sca{v}{w}=c$ and
$\sca{v}{x_{i}}\le c$ force $\sca{v}{x_{i}}=c$ whenever
$\lambda_{i}>0$. A defective witness $w\notin\X$ then lies in
$\conv(\X\cap\partial L^{+})\setminus(\X\cap\partial L^{+})$, so the
slice differs from its convex hull: it is not convex.
\end{proof}

\begin{remark}
Proposition~\ref{prop:taxonomy}(b) shows that
Conjecture~\ref{conj:revised} genuinely needs the disjunction:
defective horoballs cover reconstruction produced by asymptotic
contacts (feet escaping to infinity,
Remark~\ref{rem:escape-feet}), bitangent ones cover the
hull-collapse regime where defectiveness is blind. Whether on a
Hadamard manifold the drift branch with escaping feet always lands in
a \emph{defective} horoball is open; a robust common generalisation
would count \emph{contact ends} (limits of contact sequences in
$M\cup\bdry M$), but example (b) second part shows that even this
undercounts the defective family in the presence of escaping deep
mass. We leave the right asymptotic notion open
(Question~\ref{q:asymptotic}).
\end{remark}

\section{Sufficiency: split horoballs, and the failure of bare
bitangency}\label{sec:sufficiency}

\begin{example}[Bitangency does not imply reconstruction]
\label{ex:arc}
In the upper half-plane model of $\mathbb{H}^{2}$ let
$\X:=\{(u,1) : -1\le u\le 1\}$, a compact horocyclic arc, and let
$\gamma$ be the upward vertical ray, so that
$b_{\gamma}(x,y)=-\log y$ and $c_{\gamma}=0$. The maximal free
horoball $\Hb{\gamma}=\{y>1\}$ has contact set
$\Sb{\gamma}\cap\X=\X$: it is bitangent with a continuum of contacts.
Nevertheless $\Mx=\emptyset$: for $(a,T)\notin\X$,
\[
\cosh d\bigl((a,T),(u,1)\bigr)\;=\;1+\frac{(u-a)^{2}+(T-1)^{2}}{2T}
\]
is strictly increasing in $(u-a)^{2}$, so the nearest point of the
arc is the unique point with $u=\min(\max(a,-1),1)$. Hence
$\Rx=\emptyset\not\supseteq\Hb{\gamma}$: \emph{bitangent sufficiency
is false}, and the second disjunct of
Conjecture~\ref{conj:revised} must be the split family of
Definition~\ref{def:split}. The Euclidean shadow ($\X$ a segment in
$\R^{2}$) is handled correctly by Theorem~6 of \cite{Bia} because the
boundary slice is convex and the half-space is not defective
(Proposition~\ref{prop:taxonomy}(c)): what reconstruction detects is
a \emph{gap} in the contact set, not multiplicity.
\end{example}

\begin{definition}[Split horoball]\label{def:split}
A maximal free horoball $\Hb{\gamma}$ is \emph{split at scale
$\eta>0$} if
\[
\X\cap\{b_{\gamma}\le c_{\gamma}+\eta\}=F_{1}\cup F_{2},\qquad
\dist(F_{1},F_{2})>0,\qquad
F_{i}\cap \Sb{\gamma}\neq\emptyset\ \ (i=1,2),
\]
with $F_{1},F_{2}$ closed and nonempty.
\end{definition}

\begin{remark}\label{rem:split-vs-defective}
A split horoball is bitangent. In $\R^{n}$ a split half-space is
defective: choose contacts $x_{i}\in F_{i}\cap\partial L^{+}$ and
bisect repeatedly, always replacing the endpoint lying in the same
cluster as the midpoint; the pair distance halves, so after finitely
many steps it drops below $2\dist(F_{1},F_{2})$, at which point the
midpoint can lie in neither cluster, hence not in $\X$ --- a witness.
Thus splitness is invisible in the Euclidean statement; on a Hadamard
manifold it is the part of defectiveness that survives the collapse
of the hull (Remark~\ref{rem:collapse}).
\end{remark}

\begin{lemma}[Pure descent from a contact]\label{lem:pure}
Let $\Hb{\gamma}$ be maximal free, normalised $c_{\gamma}=0$, let
$x_{0}\in \Sb{\gamma}\cap\X$, and let $\sigma_{0}$ be the asymptotic
ray from $x_{0}$ to $\xi=\gamma(\infty)$. Then for all $t>0$:
$b_{\gamma}(\sigma_{0}(t))=-t$, $d(\sigma_{0}(t))=t$, and
$m(\sigma_{0}(t))=\{x_{0}\}$.
\end{lemma}

\begin{proof}
By (H3), $b_{\gamma}\circ\sigma_{0}$ has slope $-1$ and vanishes at
$t=0$. For $x\in\X$ the $1$-Lipschitz bound gives
$d(x,\sigma_{0}(t))\ge b_{\gamma}(x)-b_{\gamma}(\sigma_{0}(t))\ge t$,
while $d(x_{0},\sigma_{0}(t))=t$; hence $d(\sigma_{0}(t))=t$ with
$x_{0}$ a foot. If $x$ is any foot, equality
$d(x,\sigma_{0}(t))=t$ forces $b_{\gamma}(x)=0$ and unit-rate descent
of $b_{\gamma}$ along the segment from $x$ to $\sigma_{0}(t)$; with
$|\nabla b_{\gamma}|\equiv1$ (H1), the unit-speed segment satisfies
$\dot c=-\nabla b_{\gamma}(c)$, so it is an integral curve of the
locally Lipschitz field $-\nabla b_{\gamma}$ through
$\sigma_{0}(t)$, as is $\sigma_{0}$; uniqueness of integral curves
gives $x=\sigma_{0}(0)=x_{0}$ (cf.\ Step~1 of \cite[Thm.~4.1]{R1}).
\end{proof}

\begin{theorem}[Split sufficiency]\label{thm:split}
Let $M$ be Hadamard, $\X$ closed, and $\Hb{\gamma}$ split at scale
$\eta$ with clusters $F_{1},F_{2}$. Choose contacts
$x_{i}\in F_{i}\cap \Sb{\gamma}$, let $\sigma_{i}$ be the asymptotic
rays from $x_{i}$ to $\xi=\gamma(\infty)$, and set
$L_{\infty}:=\lim_{t\to\infty}d(\sigma_{1}(t),\sigma_{2}(t))$. If
$\eta>L_{\infty}/2$, then $\Hb{\gamma}\subseteq\Rx$.
\end{theorem}

\begin{proof}
Normalise $c_{\gamma}=0$ and write $b:=b_{\gamma}$; by (H3),
normalised at the contacts, $b_{\sigma_{1}}=b_{\sigma_{2}}=b$. The
function $f(t):=d(\sigma_{1}(t),\sigma_{2}(t))$ is convex and bounded
(the rays are asymptotic to the same ideal point), hence
non-increasing, with $f(t)\downarrow L_{\infty}$. Fix $t$ with
$f(t)\le2\eta$ and let $J_{t}:=[\sigma_{1}(t),\sigma_{2}(t)]$ be the
geodesic segment.

\emph{Step 1: all feet over $J_{t}$ lie in $F_{1}\cup F_{2}$.} Let
$z\in J_{t}$ and $x\in m(z)$. Convexity of $b$ along $J_{t}$ gives
$b(z)\le\max_{i}b(\sigma_{i}(t))=-t$. Since $z$ lies on a segment of
length $f(t)$, there is an $i$ with $d(z,\sigma_{i}(t))\le f(t)/2$,
whence
$d(z)\le d(z,x_{i})\le d(z,\sigma_{i}(t))+d(\sigma_{i}(t),x_{i})\le
t+f(t)/2$. The Lipschitz bound then yields
$b(x)\le b(z)+d(z,x)=b(z)+d(z)\le f(t)/2\le\eta$, so
$x\in\X\cap\{b\le\eta\}=F_{1}\cup F_{2}$.

\emph{Step 2: a medial point $z_{t}\in J_{t}$ with $d(z_{t})\ge t$.}
Set $A_{i}:=\{z\in J_{t} : m(z)\subseteq F_{i}\}$. By
Lemma~\ref{lem:pure}, $\sigma_{i}(t)\in A_{i}$, so both sets are
nonempty; they are disjoint because feet sets are nonempty and
$F_{1}\cap F_{2}=\emptyset$. Each $A_{i}$ is relatively open: if
$z_{k}\to z$ in $J_{t}$ and each $z_{k}$ has a foot
$x_{k}\in m(z_{k})\setminus F_{1}$, then $x_{k}\in F_{2}$ by Step~1;
the $x_{k}$ are bounded ($d(z_{k})\le t+f(t)/2$ with $z_{k}$ in the
compact $J_{t}$), so a subsequence converges to some $x$ with
$d(z,x)=\lim d(z_{k},x_{k})=\lim d(z_{k})=d(z)$ and
$x\in F_{2}\cap\X$ (both closed); thus $z\notin A_{1}$, and
contrapositively $A_{1}$ is relatively open; symmetrically $A_{2}$.
Since $J_{t}$ is connected, it is not the union of the two disjoint
nonempty relatively open sets $A_{1},A_{2}$: there is
$z_{t}\in J_{t}$ with feet in both clusters, hence
$\#m(z_{t})\ge2$ and $z_{t}\in\Mx$. Finally, for every $x\in\X$,
$d(x,z_{t})\ge b(x)-b(z_{t})\ge t$, so $d(z_{t})\ge t$.

\emph{Step 3: the medial balls $\B(z_{t},d(z_{t}))$ exhaust
$\Hb{\gamma}$.} Let $p\in\Hb{\gamma}$ and $\beta:=-b(p)>0$. The
quantity $d(p,\sigma_{i}(t))-t$ is non-increasing in $t$ with limit
$b_{\sigma_{i}}(p)=-\beta$, so
$d(p,\sigma_{i}(t))\le t-\beta+\varepsilon(t)$ with
$\varepsilon(t)\downarrow0$. Convexity of $d(p,\cdot)$ along
$J_{t}$ gives
$d(p,z_{t})\le\max_{i}d(p,\sigma_{i}(t))\le t-\beta+\varepsilon(t)$.
For $t$ large enough that $\varepsilon(t)<\beta$ and $f(t)\le2\eta$,
\[
d(p,z_{t})\;<\;t\;\le\;d(z_{t}),
\]
so $p\in\B(z_{t},d(z_{t}))\subseteq\Rx$.
\end{proof}

\begin{corollary}\label{cor:pinched}
If the contact rays satisfy $L_{\infty}=0$ --- in particular whenever
$M$ has sectional curvature $\le-\kappa^{2}<0$, since asymptotic
geodesics then converge --- every split horoball, at every scale
$\eta>0$, is contained in $\Rx$.
\end{corollary}

Beyond split horoballs, the following triviality records the
sufficiency of the limit objects produced by the necessity branches:

\begin{proposition}\label{prop:medial-horoball}
If a maximal free horoball $\Hb{}$ is exhausted by medial balls
(for every compact $K\subseteq\Hb{}$ there is $a\in\Mx$ with
$K\subseteq\B(a,d(a))$), then $\Hb{}\subseteq\Rx$. In particular the
drift horoballs of Theorem~\ref{thm:branch-two} satisfy
$\{b<s_{\infty}\}\subseteq\Rx$: their approximating centres are
medial by hypothesis. For the single-foot escape horoballs of
Theorem~\ref{thm:branch-one} no such conclusion is available ---
consistently with their single attained contact and with
Example~\ref{ex:arc}.
\end{proposition}

\begin{proof}
Immediate: every point of $\Hb{}$ lies in some such compact $K$.
\end{proof}

\begin{remark}[Where this leaves the conjecture]\label{rem:leaves}
Sufficiency now reads: $\Xhat\setminus\X$ always
(\cite[Thm.~3.1]{R1}); defective horoballs always
(\cite[Thm.~4.1]{R1}); split horoballs under
$\eta>L_{\infty}/2$, hence unconditionally when $L_{\infty}=0$
(Theorem~\ref{thm:split}, Corollary~\ref{cor:pinched}). Two gaps
remain. On the sufficiency side: the regime
$0<2\eta\le L_{\infty}$, where the horoball is split only at scales
below half the asymptotic separation of the contact rays (possible
when those rays bound a flat half-strip); there neither
Theorem~\ref{thm:split} nor defectiveness applies. On the necessity
side: the branches of \S\ref{sec:structure} deliver maximal free
horoballs with at least one attained contact (drift: subsequential
limits of feet; ray escape: the inflation foot), and for
reconstructible $p$ these must be upgraded to
\emph{split-or-defective} (Question~\ref{q:merge}). By
Proposition~\ref{prop:rigidity} the upgrade cannot come from a single
inflation: the second contact cluster must be injected by equidistant
motion of centres, for which the pincer of Theorem~\ref{thm:split} is
the prototype.
\end{remark}

\section{Progress on the split upgrade}\label{sec:upgrade}

Throughout this section $p\in\X^{c}\setminus\Xhat$. Normalise all
Busemann functions at a basepoint ($b_{\xi}(o)=0$); then
$(\xi,y)\mapsto b_{\xi}(y)$ is continuous on $\bdry M\times M$
\cite[II.8]{BH}, so $c_{\xi}:=\inf_{\X}b_{\xi}\in[-\infty,\infty)$ is
upper semicontinuous in $\xi$ and the \emph{depth function}
\[
D(\xi)\;:=\;c_{\xi}-b_{\xi}(p)
\]
is upper semicontinuous on the compact $\bdry M$; it is independent of
the normalisation, positive somewhere precisely because
$p\notin\Xhat$, and attains its maximum
$\delta^{\star}:=\max_{\xi}D(\xi)>0$ at some $\xi^{\star}$.

\begin{theorem}[Drift or asymptotic tangency]\label{thm:dichotomy}
Let $p\in\X^{c}\setminus\Xhat$, let $\gamma^{\star}$ be a ray
asymptotic to a depth-maximising $\xi^{\star}$, and write
$b^{\star}:=b_{\gamma^{\star}}$, $c^{\star}:=c_{\gamma^{\star}}$.
Then either $T_{\delta^{\star}/4}(p)=\infty$ (\emph{drift}), or for
all sufficiently large $t$ the inflation from $\gamma^{\star}(t)$
through any foot escapes, and in that case:
\begin{enumerate}[label=(\alph*)]
\item $d(\gamma^{\star}(t))=t+c^{\star}+o(1)$ (asymptotically exact
free radii along the deepest axis);
\item every choice of feet $x_{t}\in m(\gamma^{\star}(t))$ satisfies
$b^{\star}(x_{t})\to c^{\star}$ (asymptotic contacts of the deepest
horoball);
\item every subsequential limit $\xi'$ of the escape directions is
itself depth-maximising: $D(\xi')=\delta^{\star}$.
\end{enumerate}
\end{theorem}

\begin{proof}
For $x\in\X$ and every $t$, the non-increasing limit defining
$b^{\star}$ gives $d(x,\gamma^{\star}(t))-t\ge b^{\star}(x)\ge
c^{\star}$, so $d(\gamma^{\star}(t))\ge t+c^{\star}$; set
$\Delta_{t}:=d(\gamma^{\star}(t))-t-c^{\star}\ge0$ and
$\varepsilon(t):=d(p,\gamma^{\star}(t))-t-b^{\star}(p)\downarrow0$.
The margin of $p$ in the maximal free ball at $\gamma^{\star}(t)$ is
\[
\mu_{t}\;=\;d(\gamma^{\star}(t))-d(p,\gamma^{\star}(t))
\;=\;\Delta_{t}+\delta^{\star}-\varepsilon(t)\;>\;0
\quad\text{for large } t .
\]
If the inflations from $\gamma^{\star}(t)$ are cut for an unbounded
set of $t$, Lemma~\ref{lem:cutbranch} (with
$\varepsilon=\mu_{t}/4$) produces $a'_{t}\in\Mx$ with
$d(a'_{t})\ge d(\gamma^{\star}(t))-\mu_{t}/4\to\infty$ and
$\mu(a'_{t},p)\ge\mu_{t}/2\ge(\delta^{\star}-\varepsilon(t))/2\ge
\delta^{\star}/4$ eventually: drift. Otherwise the inflation from
$\gamma^{\star}(t)$ escapes for all large $t$;
Theorem~\ref{thm:branch-one} places $p$ at depth $\ge\mu_{t}$ in the
maximal free horoball of the escape direction $\xi_{t}$, whence
\[
\Delta_{t}+\delta^{\star}-\varepsilon(t)\;=\;\mu_{t}\;\le\;
D(\xi_{t})\;\le\;\delta^{\star},
\]
so $\Delta_{t}\le\varepsilon(t)\to0$, proving (a). For (b),
$c^{\star}\le b^{\star}(x_{t})\le
d(x_{t},\gamma^{\star}(t))-t=c^{\star}+\Delta_{t}$. For (c),
$D(\xi_{t})\ge\mu_{t}\to\delta^{\star}$ and upper semicontinuity give
$D(\xi')\ge\delta^{\star}$ for any subsequential limit, while
$\le\delta^{\star}$ by maximality.
\end{proof}

\begin{remark}
In the non-drift case the deepest horoball is asymptotically tangent
along its own axis: the obstruction to drift is \emph{exactly} an
asymptotic-contact regime. The next example shows that this regime
is, at least in the model case, the territory of defectiveness ---
the hull does \emph{not} collapse against asymptotic strands.
\end{remark}

\begin{example}[Asymptotic strand: defective but not split]
\label{ex:strand}
In the upper half-plane model of $\mathbb{H}^{2}$ let
\[
\X\;:=\;\{x_{1}\}\cup
\{(u,\,1-e^{u}) : u\le-1\},\qquad x_{1}:=(0,1),
\]
a contact point together with a strand asymptotic to the horosphere
$\{y=1\}$ from below. For the upward ray $\gamma$
($b_{\gamma}=-\log y$, $c_{\gamma}=0$), the maximal free horoball
$\Hb{\gamma}=\{y>1\}$ has the single attained contact $x_{1}$, so it
is not split (a second cluster meeting $\Sb{\gamma}$ does not exist).
It \emph{is} defective: $w:=(-10,1)$ is a trapped witness. Indeed, a
free horoball containing $w$ centred at $\infty$ is some
$\{y>c\}$ with $c<1$, which meets the strand (heights
$\to1$). A free horoball centred at $\xi\in\R$ is a Euclidean disk
tangent at $(\xi,0)$ of radius $\rho$; containing $w$ forces
$(\xi+10)^{2}+1<2\rho$, freeness against $x_{1}$ forces
$2\rho\le\xi^{2}+1$ (so $\xi<-5$), and then freeness against the
strand point $(\xi,1-e^{\xi})$ forces $2\rho\le1-e^{\xi}<1$ ---
contradicting $2\rho>(\xi+10)^{2}+1\ge1$. Hence $w$ lies in no free
horoball: $w\in(\Sb{\gamma}\cap\Xhat)\setminus\X$, and
$\Hb{\gamma}\subseteq\Rx$ by \cite[Thm.~4.1]{R1}. Defectiveness thus
covers the asymptotic-contact regime in which splitness is
unavailable, complementing Example~\ref{ex:arc}.
\end{example}

\begin{proposition}[Split upgrade on $\mathbb{H}^{2}$:
non-degenerate drift]\label{prop:h2-upgrade}
Let $M=\mathbb{H}^{2}$ and let $p$ be in the drift case, with medial
balls $B_{k}:=\B(a_{k},r_{k})\ni p$ (margins $\ge\delta$,
$r_{k}\to\infty$) as in Theorem~\ref{thm:branch-two}, whose feet
pairs $q_{k}\neq q'_{k}$ subconverge to \emph{distinct} contacts
$x^{\star}\neq x^{\star\prime}$ of the limit horoball
$\Hb{}^{\star}$. Then $\Hb{}^{\star}$ is split; consequently
(Corollary~\ref{cor:pinched}) $\Hb{}^{\star}\subseteq\Rx$, and $p$
lies in a split maximal free horoball at depth $\ge\delta$.
\end{proposition}

\begin{proof}
Work in the upper half-plane model, normalised so that the limit
direction is $\xi=\infty$, $b(x,y)=-\log y$, and the limit level
$s_{\infty}$ is the horosphere $\{y=Y\}$; then
$\X\subseteq\{y\le Y\}$ (Theorem~\ref{thm:branch-two}(a)), the
contacts are $x^{\star}=(u_{1},Y)$, $x^{\star\prime}=(u_{2},Y)$ with
$G:=u_{2}-u_{1}>0$, and $b(x^{\star})=s_{\infty}=c_{\xi}$, attained.
Hyperbolic balls are Euclidean disks, so each $B_{k}$ is a Euclidean
disk; by the exhaustion property of Theorem~\ref{thm:branch-two}(b)
the $B_{k}$ eventually contain Euclidean disks of any prescribed
radius inside $\{y>Y\}$, hence their Euclidean radii satisfy
$R_{k}\to\infty$.

\emph{Clearance.} The boundary circle of $B_{k}$ passes through
$q_{k}\to(u_{1},Y)$ and $q'_{k}\to(u_{2},Y)$, whose heights are
$\le Y$ (they lie in $\X$). Let $\ell_{k}$ be the chord through
them; $\ell_{k}\le Y$ on the segment, and the lower boundary arc lies
below $\ell_{k}$, with sagitta at horizontal distance $\ge G/4$ from
both feet at least $2\beta_{k}$, where
$\beta_{k}:=G^{2}/(128R_{k})>0$ (parabolic approximation of a circle
of radius $R_{k}\gg G$; the feet positions converge to $u_{1},u_{2}$,
so for large $k$ the middle half
$I:=[u_{1}+G/4,\,u_{2}-G/4]$ of the positions lies strictly between
them). Consequently, for large $k$, every point with position in $I$
and height in $[Y-\beta_{k},\,Y]$ lies strictly between the lower arc
and the upper arc, i.e.\ in the open disk $B_{k}$; freeness forbids
$\X$ there.

\emph{Partition.} In $b$-levels, heights $[Y-\beta_{k},Y]$ are
exactly $\{s_{\infty}\le b\le s_{\infty}+\eta_{k}\}$ with
$\eta_{k}:=-\log(1-\beta_{k}/Y)>0$, and points with $b\le
s_{\infty}+\eta_{k}$ have height $\ge Y-\beta_{k}$. Hence
$\X\cap\{b\le s_{\infty}+\eta_{k}\}$ contains no point with position
in $I$, and splits as $F_{1}\cup F_{2}$ by position ($\le u_{1}+G/4$
versus $\ge u_{2}-G/4$). Both are closed, contain the contacts
$x^{\star}$, $x^{\star\prime}$ respectively, and have positive
hyperbolic separation: any path joining them either stays below
height $H$ at horizontal cost $\ge(G/2)/H$, or climbs to height $H$
at cost $\ge2\log(H/Y)$, so
$\dist(F_{1},F_{2})\ge\min(G/(2Y),\,2\log\tfrac{G}{4Y}+2)>0$. Thus
$\Hb{}^{\star}$ is split at scale $\eta_{k}$; since
$L_{\infty}=0$ on $\mathbb{H}^{2}$, Theorem~\ref{thm:split} applies
at every scale.
\end{proof}

\begin{remark}[The remaining core]\label{rem:core}
After Theorem~\ref{thm:dichotomy} and
Proposition~\ref{prop:h2-upgrade}, what remains of
Question~\ref{q:merge} is exactly:
\begin{enumerate}[label=(\roman*)]
\item \emph{degenerate drift}: feet pairs merging
($d(q_{k},q'_{k})\to0$) or escaping to infinity --- the clearance
scale $\beta_{k}$ and the separation both degenerate with $G$;
\item \emph{the non-drift case}: upgrading the asymptotic tangency of
Theorem~\ref{thm:dichotomy} to defectiveness, for which
Example~\ref{ex:strand} is the model;
\item \emph{dimension $\ge3$}: the cleared region is a disk-shaped
shadow in the horosphere $\cong\R^{n-1}$, whose complement is
connected --- contact chains may detour around it. This is not a
mere proof-technique obstruction: \S\ref{sec:dim3} shows that for
$n\ge3$ the split-or-defective conjecture is \emph{false}, and that
splitness is a two-dimensional artefact. The remaining cases (i),
(ii) are therefore moot for Conjecture~\ref{conj:revised} as stated;
they survive only relative to the corrected notion of
\S\ref{sec:dim3}.
\end{enumerate}
\end{remark}

\section{The conjecture fails in dimension three}\label{sec:dim3}

The split refinement was forced by the collapse of the horoball hull
(Remark~\ref{rem:collapse}); the two-point set in $\mathbb{H}^{2}$
recovered reconstruction through bitangency, two \emph{separated}
contacts. We now show that in $\mathbb{H}^{3}$ a single connected
contact set already defeats Conjecture~\ref{conj:revised}: splitness
is a coincidence of dimension two.

\begin{theorem}[A counterexample in $\mathbb{H}^{3}$]\label{thm:circle}
Work in the upper half-space model
$\mathbb{H}^{3}=\{(x,y,z):z>0\}$, $\cosh d(p,q)=1+\lvert
p-q\rvert^{2}/(2z_{p}z_{q})$, and let
\[
\X\;=\;\{(\cos\theta,\sin\theta,1):\theta\in[0,2\pi)\},
\]
the unit circle on the horosphere $\{z=1\}$. Write
$R=\sqrt{x^{2}+y^{2}}$ and
$\mathcal{H}_{0}=\{R^{2}+(z-1)^{2}<1\}$ (the open Euclidean ball of
centre $(0,0,1)$ and radius $1$, a free horoball centred at the ideal
point $0$). Then
\begin{enumerate}[label=(\alph*)]
\item $\Xhat=\X$: the horoball hull collapses to the circle;
\item $\Rx=\{z>1\}\cup\mathcal{H}_{0}$, a non-empty open set; it is
exactly the union of the two maximal free horoballs whose horospheres
contain all of $\X$;
\item every maximal free horoball meets $\X$ in a connected set
(either a single point or the whole circle); hence none is split,
and by \emph{(a)} none is defective;
\item consequently the right-hand side of
Conjecture~\ref{conj:revised} is empty while $\Rx\neq\emptyset$. The
conjecture is false for $n=3$.
\end{enumerate}
\end{theorem}

\begin{proof}
The medial axis is the $z$-axis: for $a=(0,0,h)$, $h>0$,
$\cosh d(a,(\cos\theta,\sin\theta,1))=1+\frac{1+(h-1)^{2}}{2h}=
\frac{h^{2}+2}{2h}$ independently of $\theta$, so all of $\X$ is
equidistant and $a$ is medial; any point off the axis has, by strict
monotonicity of $\cosh d$ in $-2R\cos\theta$, a unique nearest point,
so is not medial.

\emph{(b)} A point $p=(x,y,z)$ lies in the medial ball
$\B(a,d(a))$ iff $\cosh d(p,a)<\cosh d(a)$, i.e.\
$1+\frac{R^{2}+(z-h)^{2}}{2zh}<\frac{h^{2}+2}{2h}$. Clearing
denominators ($2zh>0$) and cancelling $-2zh$,
\[
R^{2}+z^{2}-2z+(1-z)\,h^{2}\;<\;0 .\tag{$\ast$}
\]
If $z>1$ the coefficient $(1-z)$ is negative, so $(\ast)$ holds for
$h$ large: $\{z>1\}\subseteq\Rx$. If $z=1$, $(\ast)$ reads $R^{2}<1$.
If $z<1$, the left side is increasing in $h^{2}$, with infimum
$R^{2}+z^{2}-2z$ as $h\to0^{+}$, so $(\ast)$ is solvable iff
$R^{2}+(z-1)^{2}<1$. Since off-axis points are not medial, these
axis balls give all of $\Rx$; their union is
$\{z>1\}\cup\mathcal{H}_{0}$. The horospheres $\{z=1\}$ (ideal point
$\infty$) and $\partial\mathcal{H}_{0}$ (ideal point $0$) both contain
the entire circle.

\emph{(a)} Each $q=(x_{0},y_{0},z_{0})\notin\X$ lies in a free open
horoball: if $z_{0}>1$, in $\{z>1\}$; if $R_{0}^{2}+(z_{0}-1)^{2}<1$,
in $\mathcal{H}_{0}$; otherwise $z_{0}\le1$, $R_{0}>0$, and the
horoball centred at the ideal point $(x_{0},y_{0},0)$ --- the
Euclidean ball of centre $(x_{0},y_{0},\rho)$ and radius $\rho$ ---
is free iff $\rho\le\tfrac12((R_{0}-1)^{2}+1)$ and contains $q$ iff
$\rho>z_{0}/2$; these are compatible because
$z_{0}-1<(R_{0}-1)^{2}$ (strict, as $q\notin\X$). Hence
$M\setminus\X$ is covered by free horoballs and $\Xhat=\X$.

\emph{(c)} A maximal free horoball is centred at an ideal point
$\xi$. For $\xi=\infty$ it is $\{z>1\}$, contact set the whole
circle. For finite $\xi$ of Euclidean norm $\sigma$ it is the
Euclidean ball of centre $(\xi,\rho^{\ast})$, radius
$\rho^{\ast}=\tfrac12((\sigma-1)^{2}+1)$, and the squared Euclidean
distance from its centre to $(\cos\theta,\sin\theta,1)$ equals
$(\rho^{\ast}-1)^{2}+\sigma^{2}-2\sigma\cos\theta'+1$ ($\theta'$
measured from $\xi$), minimised at the unique $\theta'=0$ when
$\sigma>0$ --- one contact --- and constant when $\sigma=0$ --- the
whole circle ($\mathcal{H}_{0}$). In every case the contact set is
connected, so cannot be partitioned into two non-empty closed pieces
at positive distance: no maximal free horoball is split. By \emph{(a)},
$\Xhat\setminus\X=\varnothing$, so none is defective either.

\emph{(d)} The right-hand side of Conjecture~\ref{conj:revised} is
$(\Xhat\setminus\X)\cup\bigcup\{\text{split or defective}\}=
\varnothing$, whereas $\Rx=\{z>1\}\cup\mathcal{H}_{0}\neq\varnothing$.
\end{proof}

\begin{remark}[Splitness is a two-dimensional artefact]
\label{rem:artefact}
A horosphere $S_{\gamma}$ of a Hadamard manifold is intrinsically
flat, $S_{\gamma}\cong\mathbb{E}^{n-1}$. The reconstruction in
Theorem~\ref{thm:circle} is driven by the contact set $K=S_{\gamma}
\cap\X$ failing to be \emph{intrinsically convex} in $S_{\gamma}$:
the circle's flat convex hull is the disk, and the reconstructed
points sit over the disk. In dimension $n=2$ the horosphere is a
line, $S_{\gamma}\cong\mathbb{E}^{1}$, where ``convex hull strictly
larger than $K$'' is the same as ``$K$ has a gap'', i.e.\ $K$
disconnects --- which is precisely splitness. For $n-1\ge2$ a
connected set (the circle) can have strictly larger convex hull, so
splitness (a connectivity property) is strictly weaker than the
operative condition and misses the circle entirely. This is the
exact mechanism by which Conjecture~\ref{conj:revised} fails.
\end{remark}

\begin{definition}[Horospherical defect]\label{def:hdefect}
For a maximal free horoball $\Hb{\gamma}$ let $K_{\gamma}:=
\Sb{\gamma}\cap\X$ be its contact set and let
$\operatorname{hconv}_{\gamma}(K_{\gamma})$ be the closed convex hull
of $K_{\gamma}$ taken in the intrinsic flat metric of the horosphere
$\Sb{\gamma}\cong\mathbb{E}^{n-1}$. Call $\Hb{\gamma}$
\emph{horospherically defective} if
$\operatorname{hconv}_{\gamma}(K_{\gamma})\neq K_{\gamma}$.
\end{definition}

\begin{proposition}[The corrected notion passes the tests]
\label{prop:hdefect}
\leavevmode
\begin{enumerate}[label=(\alph*)]
\item \emph{(Circle.)} In Theorem~\ref{thm:circle} the horospherically
defective maximal free horoballs are exactly $\{z>1\}$ and
$\mathcal{H}_{0}$ (contact set the whole circle, flat hull the disk);
their union is $\Rx$. The single-contact horoballs are not
horospherically defective.
\item \emph{(Euclidean reduction.)} For $M=\R^{n}$ and compact $\X$,
$\Hb{}=L^{+}$ is horospherically defective iff the slice
$L\cap\X$ is non-convex iff $L$ is defective in the sense of
\cite{Bia} (Proposition~\ref{prop:taxonomy}(c)). Thus
Definition~\ref{def:hdefect} reduces to the hypothesis of Theorem~6.
\item \emph{(Planar reduction.)} For $M=\mathbb{H}^{2}$ a horoball
whose contact set is closed equals\,/\,refines splitness: a contact
set $K\subset\Sb{}\cong\mathbb{E}^{1}$ has
$\operatorname{hconv}(K)\neq K$ iff $K$ has a gap, the
on-the-horosphere shadow of the split condition
(Definition~\ref{def:split}).
\end{enumerate}
\end{proposition}

\begin{proof}
(a) Computed in the proof of Theorem~\ref{thm:circle}(c): only
$\xi\in\{0,\infty\}$ give multi-point (whole-circle) contact, with
flat hull the disk $\neq$ circle; finite $\xi\neq0$ give a single
contact, flat hull a point. The union of the two is
$\{z>1\}\cup\mathcal{H}_{0}=\Rx$ by (b) of that theorem.
(b) For compact $\X$ the face
$\cconv\X\cap L=\conv(L\cap\X)$, so the ambient defect witness
$w\in(\cconv\X\cap L)\setminus\X$ exists iff $\conv(L\cap\X)\neq
L\cap\X$; and on the hyperplane $L$ the intrinsic and ambient convex
hulls coincide. (c) On a line the closed convex hull of $K$ is the
smallest interval containing it, which differs from $K$ exactly when
$K$ omits an interior point, i.e.\ has a gap.
\end{proof}

\begin{conjecture}[Corrected necessity-and-sufficiency]
\label{conj:corrected}
Let $M$ be Hadamard and $\X$ closed. Then
\[
\Rx \;=\; (\Xhat\setminus\X)\;\cup\;
\bigcup\Bigl\{\Hb{\gamma} \;:\; \Hb{\gamma}\ \text{maximal free,
horospherically defective}\Bigr\}.
\]
\end{conjecture}

Three caveats keep this honest. First, the strand
(Example~\ref{ex:strand}) shows that contacts can be \emph{asymptotic}
--- $\X$ approaching $\Sb{\gamma}$ without meeting it --- so that
$K_{\gamma}$ undercounts; there $\Hb{\gamma}$ is defective in the
ambient sense but not horospherically defective, so the corrected
union must be read as horospherically defective \emph{or} defective,
and the genuinely correct object is presumably the flat convex hull
of the \emph{asymptotic} contact set
$\overline{\{x\in\X:b_{\gamma}(x)\to c_{\gamma}\}}$ in a suitable
horospherical compactification (Question~\ref{q:asymptotic}).
Second, for a horospherically defective $\Hb{\gamma}$ it is the
points lying \emph{over} $\operatorname{hconv}_{\gamma}(K_{\gamma})
\setminus\X$ that should be reconstructed; whether this is the whole
horoball (as for the circle and in $\R^{n}$) or a proper subregion in
general needs the inflation analysis of \S\ref{sec:inflation}
re-run with horospherical feet. Third, sufficiency
($\supseteq$) is now itself open in dimension $\ge3$: the split
sufficiency Theorem~\ref{thm:split} must be upgraded to a
horospherical-hull statement, the natural vehicle being a
\emph{wall} of medial centres swept over a triangulation of
$\operatorname{hconv}_{\gamma}(K_{\gamma})$ rather than the single
segment $J_{t}$ of the planar pincer.

\section{Sufficiency for horospherically defective horoballs}
\label{sec:hsuff}

We now prove the $\supseteq$ half of Conjecture~\ref{conj:corrected}
in real hyperbolic space $\mathbb{H}^{n}$ --- where the horospheres
are genuinely flat, $\Sb{\gamma}\cong\mathbb{E}^{n-1}$, so
Definition~\ref{def:hdefect} is unambiguous. The mechanism is a clean
duality with a classical theorem.

\begin{theorem}[Motzkin \cite{Mot}]\label{thm:motzkin}
A non-empty closed set $C\subseteq\mathbb{E}^{N}$ is convex if and only
if it is \emph{Chebyshev}: every point of $\mathbb{E}^{N}$ has a unique
nearest point in $C$. Equivalently, $C$ is non-convex iff some point
has at least two nearest points in $C$.
\end{theorem}

We use the upper half-space model
$\mathbb{H}^{n}=\{(\xi,z):\xi\in\R^{n-1},z>0\}$ with
$\cosh d(p,q)=1+\lvert p-q\rvert^{2}/(2z_{p}z_{q})$. Isometries act
transitively on horoballs and restrict to Euclidean similarities of
the flat horosphere (preserving convexity), so for any maximal free
horoball it costs nothing to assume
\[
\Hb{}=\{z>1\},\quad \Sb{}=\{z=1\}\cong\R^{n-1},\quad
\X\subseteq\{z\le1\},\quad K=\X\cap\{z=1\}.
\]
For $a=(\alpha,t)$ and a contact $k=(\kappa,1)$,
$\cosh d(a,k)=1+\frac{\lvert\alpha-\kappa\rvert^{2}+(t-1)^{2}}{2t}$ is
strictly increasing in the \emph{Euclidean} position distance
$\lvert\alpha-\kappa\rvert$; so among contacts the feet of $a$ are the
Euclidean-nearest elements of $K$, and the medial directions are the
non-Chebyshev points of $K$.

\begin{theorem}[Reconstruction on a horosphere]\label{thm:hsphere}
Let $\X$ lie on a horosphere of $\mathbb{H}^{n}$ (so $K=\X$). Then:
\begin{enumerate}[label=(\alph*)]
\item if $\X$ is not convex in $\Sb{}\cong\mathbb{E}^{n-1}$, the entire
horoball is reconstructed, $\Hb{}\subseteq\Rx$;
\item if $\X$ is convex, then $\Mx=\varnothing$ and
$\Rx=\varnothing$.
\end{enumerate}
\end{theorem}

\begin{proof}
Every point of $\X$ has height $1$, so for any $a=(\alpha,t)$ the feet
$m(a)$ are exactly the Euclidean-nearest positions of $K=\X$ to
$\alpha$, for every $t>0$.

(b) If $\X$ is convex, Motzkin gives a unique nearest position for
each $\alpha$, so every $a$ has a single foot: $\Mx=\varnothing$.

(a) If $\X$ is non-convex, Motzkin gives $\alpha^{\ast}$ with two
distinct nearest contacts $k_{1},k_{2}$ at position distance
$\rho:=\rho_{K}(\alpha^{\ast})>0$. Then $a_{t}:=(\alpha^{\ast},t)$ has
$k_{1},k_{2}\in m(a_{t})$ for every $t>0$: it is medial, with
$\cosh d(a_{t})=1+\frac{\rho^{2}+(t-1)^{2}}{2t}$. For
$p=(\pi,z)\in\Hb{}$ (so $z>1$),
$p\in\B(a_{t},d(a_{t}))$ iff $\cosh d(p,a_{t})<\cosh d(a_{t})$, which
after clearing denominators and cancelling $-2zt$ reads
\[
(1-z)\,t^{2}\;+\;\bigl(\lvert\pi-\alpha^{\ast}\rvert^{2}+z^{2}-z\rho^{2}
-z\bigr)\;<\;0 .
\]
Since $z>1$ the leading coefficient $1-z$ is negative, so the
inequality holds for all $t$ large enough: $p\in\Rx$. As $p\in\Hb{}$
was arbitrary, $\Hb{}\subseteq\Rx$.
\end{proof}

\begin{corollary}[The circle, completed]\label{cor:circle-suff}
For the unit circle of Theorem~\ref{thm:circle}, both maximal free
horoballs with whole-circle contact --- $\{z>1\}$ and $\mathcal{H}_{0}$
--- satisfy $\Hb{}\subseteq\Rx$, because the circle is non-convex on
each of their (flat) horospheres. Hence
$\Rx=\{z>1\}\cup\mathcal{H}_{0}$ equals the union of the
horospherically defective horoballs, and
Conjecture~\ref{conj:corrected} holds, as an equality, for this set.
The example that refuted Conjecture~\ref{conj:revised} thus confirms
its repair.
\end{corollary}

When $\X$ leaves the horosphere, the only obstruction is interference:
sub-horosphere mass could beat a contact for the chosen medial
direction. A uniform separation of $\X$ from the horosphere removes
it.

\begin{theorem}[Reconstruction, isolated contacts]\label{thm:hiso}
Let $\Hb{}=\{z>1\}$ be maximal free in $\mathbb{H}^{n}$ with contact
set $K$, and suppose the contacts are \emph{isolated from the bulk}:
$\X\cap\{z>1-\nu\}=K$ for some $\nu>0$. If $K$ is non-convex in
$\Sb{}\cong\mathbb{E}^{n-1}$ (i.e.\ $\Hb{}$ is horospherically
defective), then $\Hb{}\subseteq\Rx$.
\end{theorem}

\begin{proof}
Take $\alpha^{\ast}$ and $k_{1},k_{2}$, $\rho$ as before (Motzkin on
the closed set $K$). For $x=(\chi,s)\in\X$ with $s\le1-\nu$,
\[
\cosh d(a_{t},x)-\cosh d(a_{t},k_{1})\ \ge\
\frac{(t-s)^{2}}{2ts}-\frac{\rho^{2}+(t-1)^{2}}{2t}
=\frac{(1-s)(t^{2}-s)-s\rho^{2}}{2ts},
\]
and $(1-s)(t^{2}-s)-s\rho^{2}\ge\nu(t^{2}-1)-\rho^{2}>0$ once
$t>\sqrt{1+\rho^{2}/\nu}$ --- a bound independent of $x$. So beyond
this $t$ no bulk point beats the nearest contacts, $m(a_{t})\subseteq
K$ consists of the (two or more) Euclidean-nearest contacts, and
$a_{t}$ is medial. The reconstruction inequality of
Theorem~\ref{thm:hsphere}(a) is unchanged (it used only the contacts,
at height $1$), so $\Hb{}\subseteq\Rx$.
\end{proof}

\begin{remark}[What is proved, and what remains]\label{rem:hsuff}
Theorems~\ref{thm:hsphere}--\ref{thm:hiso} prove sufficiency for every
horospherically defective horoball in $\mathbb{H}^{n}$ whose contacts
are attained and isolated from the bulk; in particular the corrected
conjecture's $\supseteq$ holds, and Corollary~\ref{cor:circle-suff}
closes the circle counterexample. Three frontiers remain. \emph{(i)
Tangential bulk.} If $\X$ approaches the horosphere tangentially at a
contact (mass $(\chi_{j},1-\eta_{j})$ with
$\eta_{j}=o(\lvert\chi_{j}-\kappa\rvert)$), a fixed medial direction
can be shadowed; one expects reconstruction to survive through a
\emph{moving} medial direction, but the clean Motzkin argument no
longer applies verbatim. \emph{(ii) Curved horospheres.} Outside
$\mathbb{H}^{n}$ the horosphere need not be flat, and even
``intrinsically convex'' must be reread (the induced metric is a
Hadamard metric only under curvature pinching); Motzkin's theorem must
be replaced by its Chebyshev analogue in the relevant Alexandrov
class. \emph{(iii) Necessity.} The reverse inclusion is the
horospherical form of the upgrade (Question~\ref{q:merge}): a
reconstructible $p\notin\Xhat$ sits in a medial ball whose two feet,
followed along the structure trichotomy of \S\ref{sec:structure},
must be shown to limit onto a \emph{non-convex} contact set. By
Motzkin this is equivalent to the limit contacts failing to be
Chebyshev --- exactly the multiplicity that the two feet $q\neq q'$
encode, provided they neither merge nor are filled in by $\X$ in the
limit. This reduces necessity to a quantitative non-merging estimate,
the higher-dimensional descendant of the planar
Proposition~\ref{prop:h2-upgrade}.
\end{remark}

\section{Necessity: the Motzkin reduction and the non-merging
estimate}\label{sec:hnec}

We turn to $\subseteq$ in Conjecture~\ref{conj:corrected}. The same
duality controls it: a reconstructible point outside the hull is
reconstructed by medial centres whose feet, followed along the drift
of \S\ref{sec:structure}, must ascend to the horosphere and land on a
\emph{non-Chebyshev} --- hence by Motzkin non-convex --- contact set.
The one thing that could go wrong is that the two feet \emph{merge} in
the limit; the estimate below shows they cannot, except by collapsing
all the way onto a single contact, which is a measure-zero event we
isolate precisely.

Throughout, $\Hb{}=\{z>1\}$ is maximal free in $\mathbb{H}^{n}$,
$\X\subseteq\{z\le1\}$, $K=\X\cap\{z=1\}\neq\varnothing$, and for
$a=(\alpha,t)$, $x=(\chi,s)$ recall
$\cosh d(a,x)=1+\frac{\lvert\alpha-\chi\rvert^{2}+(t-s)^{2}}{2ts}$,
monotone in $\lvert\alpha-\chi\rvert$ at fixed heights.

\begin{proposition}[Contact feet force a defect]\label{prop:contactfeet}
If some medial centre $a$ reconstructing $p$ has two feet
$q\neq q'$ both lying in $K$, then $K$ is non-convex and $\Hb{}$ is
horospherically defective. No further hypotheses.
\end{proposition}

\begin{proof}
$q,q'$ are nearest points of $\X$ to $a=(\alpha,t)$, both in
$K\subseteq\{z=1\}$. By monotonicity at height $1$, their positions
$\kappa,\kappa'$ are Euclidean-nearest to $\alpha$ in $K$, at a common
distance; since $\kappa\neq\kappa'$, the point $\alpha$ has two nearest
points in $K$. So $K$ is not Chebyshev, hence by
Theorem~\ref{thm:motzkin} not convex.
\end{proof}

\begin{lemma}[Ascent]\label{lem:ascent}
Let $a_{k}=(\alpha_{k},t_{k})$ be medial with $t_{k}\to\infty$ and
$\alpha_{k}\to\alpha_{\infty}$, and let $q_{k}=(\beta_{k},s_{k})\in
m(a_{k})$ have bounded positions $\beta_{k}$. Then
$1-s_{k}=O(t_{k}^{-2})$; in particular every subsequential limit of
$q_{k}$ lies in $K$, and $h_{k}:=(1-s_{k})t_{k}^{2}$ is bounded.
\end{lemma}

\begin{proof}
Fix a contact $k_{0}=(\kappa_{0},1)\in K$. As $q_{k}$ is a nearest
point, $\cosh d(a_{k},q_{k})\le\cosh d(a_{k},k_{0})$, i.e.
\[
\frac{\lvert\alpha_{k}-\beta_{k}\rvert^{2}+(t_{k}-s_{k})^{2}}{s_{k}}
\;\le\;\lvert\alpha_{k}-\kappa_{0}\rvert^{2}+(t_{k}-1)^{2}.
\]
Dropping the non-negative position term and using
$\frac{(t_{k}-s_{k})^{2}}{s_{k}}-(t_{k}-1)^{2}=
\frac{(1-s_{k})(t_{k}^{2}-s_{k})}{s_{k}}$ gives
$\frac{(1-s_{k})(t_{k}^{2}-s_{k})}{s_{k}}\le
\lvert\alpha_{k}-\kappa_{0}\rvert^{2}\le C$, whence
$(1-s_{k})t_{k}^{2}\le C'$ for large $k$. Closedness of $\X$ then
places limits of $q_{k}$ in $\X\cap\{z=1\}=K$.
\end{proof}

\begin{theorem}[Non-merging estimate]\label{thm:nonmerge}
Let $a_{k}=(\alpha_{k},t_{k})$ be medial centres of a drift
reconstructing $p\in\Hb{}$ ($t_{k}\to\infty$, $\alpha_{k}\to
\alpha_{\infty}$, margins bounded below), with distinct feet
$q_{k}\neq q_{k}'$ of bounded position. Passing to a subsequence, let
$q_{k}\to q_{\ast}=(\kappa_{\ast},1)$ and $q_{k}'\to
q_{\ast}'=(\kappa_{\ast}',1)$ in $K$ (Lemma~\ref{lem:ascent}). Then:
\begin{enumerate}[label=(\alph*)]
\item the heights ascend sharply, $h_{\ast}=\lim(1-s_{k})t_{k}^{2}=0$
and likewise for $q_{k}'$; both $\kappa_{\ast},\kappa_{\ast}'$ are
Euclidean-nearest to $\alpha_{\infty}$ in $K$;
\item \emph{(estimate)} either $\kappa_{\ast}=\kappa_{\ast}'$
(\emph{collapse}), or $\kappa_{\ast}\neq\kappa_{\ast}'$ and then
\[
d(q_{k},q_{k}')\ \longrightarrow\
\operatorname{arccosh}\!\Bigl(1+\tfrac12\lvert\kappa_{\ast}-
\kappa_{\ast}'\rvert^{2}\Bigr)\;>\;0 ;
\]
there is no intermediate decay rate;
\item in the non-collapse case $K$ has two distinct nearest points to
$\alpha_{\infty}$, so by Motzkin it is non-convex: $\Hb{}$ is
horospherically defective and contains $p$. Necessity holds.
\end{enumerate}
\end{theorem}

\begin{proof}
Write $h_{k}=(1-s_{k})t_{k}^{2}$, $h_{k}'=(1-s_{k}')t_{k}^{2}$, bounded
by Lemma~\ref{lem:ascent}; pass to limits $h_{\ast},h_{\ast}'\ge0$.
For any contact $(\kappa,1)$ the nearest-point inequality
$\cosh d(a_{k},q_{k})\le\cosh d(a_{k},(\kappa,1))$, multiplied by
$2t_{k}$ and rearranged with
$\frac{(t_{k}-s_{k})^{2}}{s_{k}}-(t_{k}-1)^{2}=
\frac{(1-s_{k})(t_{k}^{2}-s_{k})}{s_{k}}\to h_{\ast}$, passes to the
limit as
\[
\lvert\alpha_{\infty}-\kappa_{\ast}\rvert^{2}+h_{\ast}\ \le\
\lvert\alpha_{\infty}-\kappa\rvert^{2}\qquad(\forall\,\kappa\in K),
\tag{$\dagger$}
\]
and symmetrically
$\lvert\alpha_{\infty}-\kappa_{\ast}'\rvert^{2}+h_{\ast}'\le
\lvert\alpha_{\infty}-\kappa\rvert^{2}$. Putting $\kappa=
\kappa_{\ast}'$ in the first and $\kappa=\kappa_{\ast}$ in the second
and adding yields $h_{\ast}+h_{\ast}'\le0$, so $h_{\ast}=h_{\ast}'=0$:
this proves (a), and ($\dagger$) becomes
$\lvert\alpha_{\infty}-\kappa_{\ast}\rvert=
\lvert\alpha_{\infty}-\kappa_{\ast}'\rvert=
\min_{\kappa\in K}\lvert\alpha_{\infty}-\kappa\rvert$.

For (b), $d(q_{k},q_{k}')\to d(q_{\ast},q_{\ast}')$, and for two points
at height $1$,
$\cosh d(q_{\ast},q_{\ast}')=1+\tfrac12\lvert\kappa_{\ast}-
\kappa_{\ast}'\rvert^{2}$. If $\kappa_{\ast}\neq\kappa_{\ast}'$ this is
$>1$, the stated bound; the absence of intermediate rates is the
dichotomy ``equal or not''. For (c), distinct equidistant nearest
points show $K$ non-Chebyshev; Motzkin gives non-convexity, and
$p$ lies in $\Hb{}=\{z>1\}$ by construction.
\end{proof}

\begin{remark}[The collapse residue]\label{rem:collapse-residue}
The estimate leaves exactly one way for necessity to be threatened:
\emph{collapse}, $q_{k},q_{k}'\to$ a single contact $\kappa_{\ast}$.
Three remarks bound its reach. First, collapse cannot occur over a
\emph{smooth} contact: if near $\kappa_{\ast}$ the set $\X$ is a graph
$z=1-\phi$ with $\phi\ge0$ strictly convex (a rounded touching), the
high-centre feet minimise $\lvert\alpha-\chi\rvert^{2}+t^{2}\phi(\chi)$
which is strictly convex, giving a \emph{unique} foot --- $a_{k}$ is
then not medial. Collapse therefore forces a non-smooth (cusp or
concave-corner) contact, where $K$ is already non-convex, so
Proposition~\ref{prop:contactfeet}/Motzkin still apply provided the
feet are eventually \emph{in} $K$. Second, the genuinely residual case
is collapse with feet strictly below the horosphere hugging an
isolated contact; there the reconstructed points lie within $o(1)$ of
$\kappa_{\ast}$, and whether they fall in $\Xhat$ (first term, nothing
to prove) or expose an ambient defect is exactly the local-hull
content of the asymptotic-contact Question~\ref{q:asymptotic}. Third,
when contacts escape to infinity ($\beta_{k}$ unbounded) the limit is
asymptotic and again governed by Question~\ref{q:asymptotic} (the
strand, Example~\ref{ex:strand}). Thus: \emph{necessity is proved in
$\mathbb{H}^{n}$ for every drift whose feet stay bounded and do not
collapse onto one contact; the remainder reduces, with no loss, to the
single asymptotic/local-hull question already isolated as open.}
\end{remark}

\section{Summary and open problems}\label{sec:summary}

\begin{center}
\begin{tabular}{lll}
\toprule
Result & Hypothesis & Status \\
\midrule
$\Xhat\setminus\X\subseteq\Rx$ & complete & proved \cite[Thm.~3.1]{R1} \\
defective $\Hb{}\subseteq\Rx$ & Hadamard & proved \cite[Thm.~4.1]{R1} \\
No escape from medial centres (rigidity) &
Hadamard & proved (Prop.~\ref{prop:rigidity}) \\
Drift $\Rightarrow$ maximal $\Hb{}\ni p$, contacts captured &
Hadamard & proved (Thm.~\ref{thm:branch-two}) \\
Trichotomy hull/drift/ray-escape, no residual case &
Hadamard & proved (Thm.~\ref{thm:structure}) \\
$\Rx\setminus\Xhat\subseteq$ drift & $\R^{n}$ & proved
(Prop.~\ref{prop:euclidean}) \\
Drift or asymptotic tangency & Hadamard & proved
(Thm.~\ref{thm:dichotomy}) \\
Split upgrade, non-degenerate drift & $\mathbb{H}^{2}$ & proved
(Prop.~\ref{prop:h2-upgrade}) \\
Defective $\setminus$ split $\neq\emptyset$ (strand) &
$\mathbb{H}^{2}$ & example (Ex.~\ref{ex:strand}) \\
Split sufficiency, $\eta>L_{\infty}/2$ & $\mathbb{H}^{2}$/pinched &
proved (Thm.~\ref{thm:split}, Cor.~\ref{cor:pinched}) \\
\midrule
\textbf{Split-or-defective conjecture} & $n\ge3$ &
\textbf{false} (Thm.~\ref{thm:circle}) \\
Splitness is a $2$-dim artefact & --- & Rem.~\ref{rem:artefact} \\
Horospherical defect: circle test & $\mathbb{H}^{3}$ & proved
(Prop.~\ref{prop:hdefect}a) \\
Horospherical defect $=$ defective & $\R^{n}$ & proved
(Prop.~\ref{prop:hdefect}b) \\
Horospherical defect $\approx$ split & $\mathbb{H}^{2}$ & proved
(Prop.~\ref{prop:hdefect}c) \\
Corrected conjecture & Hadamard & open
(Conj.~\ref{conj:corrected}) \\
\midrule
Horosph.\ defective $\Hb{}\subseteq\Rx$, $\X$ on horosphere &
$\mathbb{H}^{n}$ & proved (Thm.~\ref{thm:hsphere}) \\
$\X$ convex on horosphere $\Rightarrow\Rx=\varnothing$ &
$\mathbb{H}^{n}$ & proved (Thm.~\ref{thm:hsphere}) \\
Horosph.\ defective $\Hb{}\subseteq\Rx$, isolated contacts &
$\mathbb{H}^{n}$ & proved (Thm.~\ref{thm:hiso}) \\
Circle: $\Rx=$ union of horosph.\ defective &
$\mathbb{H}^{3}$ & proved (Cor.~\ref{cor:circle-suff}) \\
Sufficiency, tangential bulk / curved horosph. &
Hadamard & open (Rem.~\ref{rem:hsuff}) \\
\midrule
Contact feet $\Rightarrow$ defect (Motzkin) & $\mathbb{H}^{n}$ &
proved (Prop.~\ref{prop:contactfeet}) \\
Ascent of bounded feet to $K$ & $\mathbb{H}^{n}$ & proved
(Lem.~\ref{lem:ascent}) \\
Non-merging estimate $\Rightarrow$ necessity & $\mathbb{H}^{n}$ &
proved (Thm.~\ref{thm:nonmerge}) \\
Necessity, collapse / asymptotic residue & $\mathbb{H}^{n}$ &
open $=$ Q.~\ref{q:asymptotic} (Rem.~\ref{rem:collapse-residue}) \\
\bottomrule
\end{tabular}
\end{center}

\begin{question}[The necessity gap, horospherical form]
\label{q:merge}
Theorem~\ref{thm:circle} retires the split-or-defective form of the
question; the live version is whether
Conjecture~\ref{conj:corrected} holds: for $p\in\Rx\setminus\Xhat$,
must some maximal free horoball containing $p$ be horospherically
defective (Definition~\ref{def:hdefect})? By
Proposition~\ref{prop:rigidity} the upgrade cannot come from a single
inflation: the medial witness of $p$ injects its further contacts
through equidistant motion of centres, as in the pincer of
Theorem~\ref{thm:split}. The concrete target is an
\emph{equidistant-flow theorem}: starting from $a\in\Mx$ and
$p\in\B(a,d(a))$, grow free balls with centres on the equidistant
locus of the contact set, and show the family drifts to a horoball
whose horospherical contact hull contains the shadow of $p$. For
$n=2$ this is the split upgrade of \S\ref{sec:upgrade}
(Remark~\ref{rem:core} delimits its open core); for $n\ge3$ the
single segment $J_{t}$ of the planar pincer must become a swept
$(n-1)$-dimensional wall over
$\operatorname{hconv}_{\gamma}(K_{\gamma})$.
\end{question}

\begin{question}\label{q:asymptotic}
Formulate the correct asymptotic notion of horospherical contact
hull --- the flat convex hull in $\Sb{\gamma}\cong\mathbb{E}^{n-1}$ of
the \emph{asymptotic} contact set
$\overline{\{x\in\X:b_{\gamma}(x)\to c_{\gamma}\}}$, in a horospherical
compactification --- so that Conjecture~\ref{conj:corrected} both
covers the asymptotic strand (Example~\ref{ex:strand},
Proposition~\ref{prop:taxonomy}(b)), where contacts recede to the
ideal boundary, and reduces to Definition~\ref{def:hdefect} when
contacts are attained. Then identify, for a horospherically
defective horoball, the exact reconstructed subregion (the second
caveat after Conjecture~\ref{conj:corrected}), and prove the
resulting equality.
\end{question}

\begin{question}\label{q:positive-curv}
Theorem~\ref{thm:branch-one} uses only global minimality of the
extended foot-geodesic. On a complete, non-Hadamard manifold the
statement survives whenever the inflation ray is globally minimizing;
the flat cylinder (where bitangency without multiplicity is possible,
$\X=S^{1}\times\{0\}$) shows that bitangency must then be read as
``two minimizing geodesics''. Determine the sharp generality.
\end{question}

\begin{thebibliography}{9}

\bibitem{Bia} A.~Bia\l{}o\.zyt, \emph{Sets reconstructible with medial
axis}, preprint, AGH University of Krak\'ow, 2024.

\bibitem{R1} \emph{Sets reconstructible with medial axis: a
generalisation to Riemannian manifolds}, companion analysis notes,
June 2026 (file \texttt{medial\_axis\_riemannian\_1.pdf}).

\bibitem{Notes} \emph{Formalization notes on Theorem~6 of
\cite{Bia}}, June 2026 (files
\texttt{medial\_formalization\_notes.tex},
\texttt{medial\_formalization\_plan.md}); Lean~4 formalization on
branch \texttt{medial-axis}, files
\texttt{Mathlib/Analysis/Convex/MedialAxis*.lean}.

\bibitem{BH} M.~Bridson and A.~Haefliger, \emph{Metric Spaces of
Non-Positive Curvature}, Grundlehren 319, Springer, 1999.

\bibitem{Alb} P.~Albano, \emph{On the cut locus of closed sets},
Nonlinear Anal.\ \textbf{125} (2015), 398--405.

\bibitem{Wol} F.-E.~Wolter, \emph{Cut loci in bordered and unbordered
Riemannian manifolds}, PhD thesis, TU Berlin, 1985.

\bibitem{Mot} T.~Motzkin, \emph{Sur quelques propri\'et\'es
caract\'eristiques des ensembles born\'es non convexes}, Atti Accad.\
Naz.\ Lincei \textbf{21} (1935), 773--779. (Closed Chebyshev sets in
$\mathbb{E}^{N}$ are convex.)

\end{thebibliography}

\end{document}
